\chapter{Metodologia}
\label{chap:metodologia}

\section{Descrição da Ferramenta MNE}\label{sec:exemplo-de-algoritmos-e-figuras}

    O MNE-Python (cuja sigla representa \textit{Magnetoencephalography and
    Electroencephalography} em Python) é um pacote de código aberto em Python
    projetado para a exploração, visualização e análise de dados
    neurofisiológicos humanos. Tornou-se um padrão na pesquisa em neurociência
    para o processamento de sinais cerebrais devido à sua flexibilidade e
    abrangência.

    O MNE-Python suporta diversos tipos de dados neurofisiológicos, incluindo
    EEG (\textit{Electroencephalography}), MEG
    (\textit{Magnetoencephalography}), sEEG (EEG intracraniano), ECoG
    (\textit{Electrocorticography}) e fNIRS (\textit{Functional Near-Infrared
    Spectroscopy}). A biblioteca oferece um ecossistema completo de ferramentas
    que cobrem todas as etapas da pipeline de análise, desde a leitura e
    importação de dados em diversos formatos até a etapa de pré-processamento,
    análise e visualização dos resultados.

    Entre as principais funcionalidades de pré-processamento, o MNE-Python
    oferece ferramentas robustas para filtragem digital, detecção de eventos e
    remoção de artefatos. A remoção de artefatos pode ser realizada por meio de
    ICA (\textit{Independent Component Analysis}) ou SSP (\textit{Signal Space
    Projection}), técnicas amplamente utilizadas para eliminar artefatos
    oculares, cardíacos e musculares dos sinais.

    O toolkit possui capacidades de análise no domínio do tempo e no domínio
    tempo-frequência, utilizando transformadas como Fourier e wavelets de
    Morlet. Além disso, o MNE-Python implementa soluções para o problema
    inverso, permitindo a estimação de fontes neurais a partir das medições de
    sensores, utilizando métodos como \textit{beamforming}. A biblioteca também
    integra ferramentas de conectividade, estatística e aprendizado de máquina,
    tornando-a uma solução abrangente para pesquisa em neurociência.

    A estrutura de dados central do MNE-Python é organizada em três classes
    principais: \textit{Raw} para dados contínuos, \textit{Epochs} para
    segmentos de interesse delimitados por eventos, e \textit{Evoked} para
    médias de épocas associadas a condições experimentais específicas.
